\vsssub
\subsubsection{点输出后处理器} \label{sec:ww3outp}
\vsssub

\proddefH{ww3\_outp}{w3outp}{ww3\_outp.ftn}
\proddeff{输入}{ww3\_outp.inp}{传统配置文件。}{10} (App.~\ref{sec:config161})
\proddefa{mod\_def.ww3}{模式定义文件。}{20}
\proddefa{out\_pnt.ww3}{原始点输出数据。}{20}
\proddefa{NC\_globatt.inp}{附加全局属性。}{994}
\proddeff{输出}{standard out}{程序的格式化输出。}{6}
\proddefa{tab{\sl{nn}}.ww3 \opt}{平均参数表，其中
{\file{\sl{nn}}} 为两位整数。}{\sl nn}
\proddefa{\ldots \opt}{传输文件。}{user}

\vspace{\baselineskip} 
\vspace{\baselineskip} 
\noindent 
在 \ws{} 早期版本中，谱公报使用由 {\F itype = 1} 和 {\F otype = 3} 生成的谱数据传输文件以及 {\file w3split} 程序（见 section~\ref{sec:install}）生成。这是一段即将淘汰的代码，此处仅为向后兼容而提供。该程序从标准输入读取以下五条记录（不允许注释行）：

\begin{list}{$\bullet$}{\itemsep 0mm \parsep 0mm}
\item 输出位置名称。
\item 表中使用的运行标识符。
\item 输入文件名称。
\item 标识 UNFORMATTED 输入文件的逻辑量。
\item 输出文件名称。
\end{list}

\noindent
以上所有字符串均以自由格式按字符读取，因此需要用引号括起。

\pb
